\section{Perturbation model-agnostic methods}

In this section, we present a set of perturbation method used in black-box 
models explanation. A perturbation method can be defined by the 
generation of perturbed samples (ex: noised features in model representation
or perturbed input tensor) at the input of the black-box model. The objective 
is to compare the model behavior (the prediction) between a set of different 
perturbed inputs and the original one. Then, we can determine 
which one is more important for the prediction. It allows the explanation model to 
detail why the model took a specific decision and why. The following methods register in this 
type. The first one, LIME (Local Interpretable Model-Agnostic Explanation) focus on a linear approximation 
of local model behavior. Secondly, the Anchors method provides an amelioration of LIME by introducing the concept 
of IF-THEN rules to provide a better user-experience. Finally, we will briefly talk about add-ons of those methods
in specific cases. 

\begin{remark}[Notation]
    Let's introduce some notations to identify which objects we will talk about. 
    We denote with $f$ a black-box model we want to explain, then $f : X \longrightarrow Y$ 
    where $X$ is the set of instances $x \in X$ and $Y$ the different possible predictions $f(x)$. 
    In this section, we will talk about linear explanation model, denoted $g$. 
    To identify which input representation we use, we denote by $x'$ the 
    interpretable representation of an input $x$ in the black-box model $f$. 
    Then, if we construct a set of perturbed interpreted samples $ \mathcal{Z}$ 
    from $x'$, this set contains perturbed version of $x'$ called $z'$. Finally, 
    $ \Omega(g)$ express the complexity of the model $g$, generally defined by the 
    number of non-zeros weights in $g$. 
\end{remark}

\subsection{LIME (Local Interpretable Model-Agnostic Explainer)}

% Model introduction 
LIME founded the principe of perturbation method for classifiers explanation \cite{ribeiro_why_2016}. 
It provides a local explanation of a black-box model behavior 
$ f : \R^d \longrightarrow \R $ by training a simple explanation 
model $g$ over perturbed samples. This method 
is \emph{model-agnostic} because it doesn't depend on the black-box model 
architecture. $f$ is considered as an oracle used to train $g$.

% LIME technical aspects
Formally, for an original interpretable representation $x' \in \{0, 1\}^{d'}$ of a
sample $x \in \R^d$, it generates a set of perturbed samples $z'$ which are drawn
by selecting uniformly some $x'$ non-zero components. This random sampling is weighted by 
$\pi_x$ which defines a locality around $x'$ denoted $ \mathcal{Z}$. In practical, 
$ \pi_x$ is an exponential kernel applied to a distance $D$. 
To find the optimal model $g$, we can define a loss function $ \mathcal{L}(f,g, \pi_x)$ telling how 
unfaithful $g$ is in approximating $f$ on $ \mathcal{Z}$. The objective is then to 
minimize $ \mathcal{L}$ and having $ \Omega(g)$ low enough for human comprehension. 
The LIME explanation is then produced by: 
\begin{equation}
    \xi(x) = \arg \min_{g \in G} \mathcal{L}(f,g, \pi_x) + \Omega(g)
\end{equation}
with $G$ the set of available explanation models \cite{ribeiro_why_2016}. Using $ \mathcal{L}$ and $ \Omega(g)$ 
to find the best model $g$ ensures \emph{interpretability} and \emph{local fidelity}
on the black-box model explanation. Finally, the local explanation of the model 
is provided by the weights $w_g$ trained with $g(z') = w_g \cdot z'$.

% Global interpretation 
As stated at beginning, LIME provides a \emph{local} explanation of $f$. To ensure 
a global explanation, the paper proposes a Sub-modular Picks LIME (SP-LIME) 
which selects a set local and non-redundant LIME explanations to be given to the user \cite{ribeiro_why_2016}.
To choose those explanations, the paper introduces the concept of coverage of a set 
of instances which defines the importance of those instances. 

\begin{definition}[Coverage]
    Let $g$ be an explanation model and $ g_i = \xi(x_i)$ the explanation 
    of the instance $x_i$, we construct the $|X| \times d'$ explanation matrix $ \mathcal{W}$ 
    with $ \mathcal{W}_{i,j} = | w_{g_{i,j}} |$. Then, $I_j$ is the $j$-column of 
    $W$ and denotes the global importance of that component in the explanation space. 
    Then, the coverage of a set of instances $V$ is defined as the total importance 
    of the features that appear at least one instance in $V$: 
    \begin{equation}
        c(V, \mathcal{W}, I) := \sum_{j=1}^{d'} \mathbbm{1}_{\{ \exists i \in V, \mathcal{W}_{i,j} > 0\}} I_j  
    \end{equation}
\end{definition}

Finally, from a maximum number of global explanations $B$ set by the user, the 
pick problem is defined by:
    \begin{equation} \label{eq:}
        Pick( \mathcal{W},I) = \underset{V, |V| \leqslant B}{\arg \max} \; c(V, \mathcal{W}, I) 
    \end{equation}
We want to select the highest coverage and non-redundant explanations to provide a global 
one. 


% Computing and limits
Unfortunately, this system has limits. First, in practice, the pick problem is intractable, 
so a greedy algorithm is used to compute it. 
Secondly, LIME algorithm doesn't define clear boundaries of the explanation instances. Indeed, 
all black-box models cannot be locally approximated by a linear one \cite{ribeiro_nothing_2016}. 
Finally, the explanation results depend on the LIME parameters such as $B$, the distance metric 
$D$ and the kernel $ \pi_x$. Changing one of them can change the interpretation and 
undermining the reliability of this method \cite{wilking_fooling_2022}.

\subsection{Anchors}

Anchors is also a model-agnostic explanation method based on LIME system but enhanced in order
to provide local sufficient conditions for prediction. Anchors are defined by short and 
disjoint rules which are preferred by users for explanation \cite{ribeiro_anchors_2018}. It can be applied 
on different set of domains (tabular, text, images, etc...) for different 
tasks (classification, prediction, text generation, etc...).

\begin{definition}[Anchors]
    Formally, Anchors are built on a set of predicates $a$ called a "rules". We can 
    also see it as a function of instances $x \in X$ :
    $$ A : 
        \begin{cases}
            X \longrightarrow \{0, 1\} \\ 
            x \longmapsto 
            \begin{cases}
                1 \text{ if all rules are true for } x \\
                0 \text{ otherwise}
            \end{cases}
        \end{cases}
    $$
    to link a set of rules $A$ on a prediction $f(x)$ (where $f : X \longrightarrow Y$ 
    is the black-box model). Finally, $A$ is an Anchor of $x$ if $A(x) = 1$ 
    and $A$ is a sufficient condition for the prediction $y = f(x)$ :
        \begin{equation}
            \mathbb{E}_{\mathcal{D}(z \mid A)}[1_{f(x)=f(z)}] \geqslant \tau
        \end{equation}
    with $\tau$ the desired level of precision of the Anchor on $x$ \cite{ribeiro_anchors_2018}.
    Here, $ \mathcal{D}(x \mid A)$ represents the distribution of neighbors of $x$ 
    which satisfy the Anchor $A$. 
\end{definition}

Computing Anchors is the difficult part of the process. This is a combinatorial 
optimization problem if we decide to choose anchors which satisfy a precision 
constraint with high probability:
    $$	\mathbb{P} (prec(A) \geqslant \tau) \geqslant 1 - \delta$$
with $ prec(A) = \mathbb{E}_{\mathcal{D}(z \mid A)}(1_{f(x)=f(z)})$ called the Anchors 
"precision". Among those anchors, we choose the largest coverage ones, where coverage is :
\begin{equation}
    cov(A) = \mathbb{E}_{\mathcal{D}(z)} (A(z))
\end{equation}
Finally, the optimization problem can be described by :
\begin{equation}
    \max_{\mathbb{P} (prec(A) \geqslant \tau) \geqslant 1 - \delta} cov(A)
\end{equation}

Because the number of all possible Anchor is exponential, we can perform an 
exhaustive search. That's why the article introduced the Bottom-Up approach:
\begin{enumerate}
    \item \emph{Initialization}: The anchor $A$ is initialized as an empty rule (applying to all instances).
    \item \emph{Candidate Generation}: In each iteration, a set of candidate rules $\mathcal{A}$ is generated by extending $A$ with one additional feature predicate $\{a_i\}$: $\mathcal{A} = \{A \wedge a_i, A \wedge a_{i+1}, A \wedge a_{i+2}, ...\}$
    \item \emph{Selection via Multi-Armed Bandit (MAB)}: To identify the best 
    candidate without an exhaustive number of calls to $f$, the problem is 
    formulated as a pure-exploration multi-armed bandit problem. The \emph{KL-LUCB} 
    algorithm is used to find the rule with the highest precision using a minimal 
    set of samples \cite{ribeiro_anchors_2018}.
    \item \emph{Termination}: The process terminates when a rule meets the statistical precision constraint: $\mathbb{P} (prec(A) \geqslant \tau) \geqslant 1 - \delta$.
\end{enumerate}

To address the limitations of a purely greedy search, the algorithm extends the 
construction to a \emph{Beam Search} of size $B$.
It maintains a set of $B$ current candidates instead of a single one.
Among the multiple anchors identified that satisfy the precision threshold $\tau$, 
the algorithm outputs the one with the highest coverage $cov(A)$, effectively optimizing the original problem.

To provide a global understanding of the model, the article adapts a \emph{Submodular Pick} (SP) procedure. 
The goal is to select a set of $K$ anchors that are non-redundant and cover as 
many instances in the validation set similar to the approach used in LIME \cite{ribeiro_anchors_2018,ribeiro_why_2016}.